对固定索引，查询先变换为 $q'=\Pi q$，再路由到选中的倒排列表，使用主码筛选，仅对保留候选做精排。最终选择器返回 $k$ 个数据库 ID 及计算所得分数。\Cref{fig:5-query-path}区分运行时的三个选择决策；ground-truth 排序属于外部评价对象。

需要分清四种量：
\begin{align}
 s_i(q)&=\|\mu_i\|_2^2-2\langle q',\mu_i\rangle,
                   &&\text{centroid routing score},\nonumber\\
 \primarydist(q,x)&=\|q'-\widehat p(x)\|_2^2,
                   &&\text{primary score},\nonumber\\
 \refineddist(q,x)&=\|q'-\widehat p(x)-\widehat r(x)\|_2^2
                       +\|q'\|_2^2,
                   &&\text{refined ranking score},\nonumber\\
 \rawdist(q,x)&=\|q-x\|_2^2,
                   &&\text{raw-space distance}.
 \label{eq:5-ranking-objects}
\end{align}
其中 $\mu_i$ 是变换坐标中的 IVF 中心，$\widehat p(x)$ 为主码重建，$\widehat r(x)$ 为解码后的残差。运行时精排分数包含一个仅与查询有关的常数；\Cref{sec:5-6}依据实际修正项推导这一点。对固定查询，它保持重建距离的排序，但精排分数和主码分数都不等于原始向量距离。

\begin{figure}[tbp]
\centering
% Conceptual dataflow; no measured values or GPU timing claims.
\begin{tikzpicture}[font=\small,
 box/.style={draw=ThesisNavy!65,rounded corners=2pt,align=center,
  text width=3.7cm,minimum height=3cm,inner sep=5pt},
 arr/.style={-{Latex[length=2mm]},draw=ThesisNavy!80,thick}]
\node[box,fill=ThesisNavy!5] (route) at (0,0)
 {$q\to q'=\Pi q$\\[3pt]
  Centroid scores $s_i$\\Select $n_{\rm list}\to n_{\rm probe}$ lists\\Visited universe $U$};
\node[box,fill=ThesisNavy!5] (primary) at (4.9,0)
 {Build LUTs\\Read primary codes\\[3pt]
  Score $n_c$ candidates\\Retain primary top-$c$};
\node[box,fill=ThesisNavy!5] (refine) at (9.8,0)
 {Read residual codes\\Compute $\refineddist$\\[3pt]
  Select $c\to k$\\Return IDs and scores};
\draw[arr] (route) -- node[above,font=\footnotesize] {$n_c$} (primary);
\draw[arr] (primary) -- node[above,font=\footnotesize] {$c$} (refine);
\node[draw=ThesisNavy!40,rounded corners=2pt,align=center,
 text width=12.6cm,inner sep=6pt] at (4.9,-2.7)
 {Evaluation only: compare returned IDs with raw-space ground truth under
  $\rawdist$\\Ground truth is outside the compressed search decisions.};
\end{tikzpicture}

\caption{查询执行及其独立选择边界。路由选择列表，主码筛选选择向量候选，最终选择对精排后的保留候选排序。原空间 ground truth 用于评价，不作为这些搜索决策的输入。}
\label{fig:5-query-path}
\end{figure}

路由计算中心分数并选择 $n_{\rm probe}$ 个列表（\Cref{sec:5-7}）。拼接这些列表得到大小为 $n_c$ 的已访问集合 $U$。精确探针原语只是按已计算路由分数选择所需中心，不证明 $U$ 包含全部原空间近邻。前缀偏移还给出扫描与调度所需的实际候选数。

主码筛选构造每查询 LUT，读取已访问向量的主码，再按 $\primarydist$ 保留至多 $c=\max\{k,\lceil\alpha k\rceil\}$ 个候选。固定表分组时，分数计算成本为 $O(n_cM)$，此外还有局部选择和合并。\Cref{sec:5-2,sec:5-3}讨论表与编码访问，\Cref{sec:5-4,sec:5-5}建立选择不变量，并分析增加并行度的成本。

残差评分只对有效保留候选读取残差码，进行 $O(cd)$ 坐标工作，并向最终选择器提供 $\refineddist$（\Cref{sec:5-6}）。这一路径使用压缩重建替代搜索时的原始向量读取。实现中的最终 repeated-minimum 选择器为每个输出扫描 $c$ 个槽位，共输出 $k$ 个，另有 block 归约工作。

确定性排序 key 为各类分数附加并列规则：中心选择使用列表 ID，完整 key 主码选择使用物理行位置，最终选择使用数据库 ID。旧的查询级主码路径具有不同的截断并列规则，见\Cref{sec:5-4,sec:5-5}的限定。因此，分数保持、保留集合一致、有序输出相等及原空间 recall 是不同主张。第七章在各自边界上检查它们，而不把单个 recall 不变当作整条流水线的证明。
